%==============================================================================
%  ELLIPTIC CURVE CRYPTOGRAPHY -- Lecture Notes for MPhil/PhD students
%  Compile with:  pdflatex main.tex  (run twice)   or   latexmk -pdf main.tex
%  Requires: a standard TeX Live installation (pgfplots, tcolorbox, listings,
%  algorithmicx, cleveref, hyperref are all included in TeX Live "full").
%==============================================================================
\documentclass[11pt,a4paper,oneside]{book}

\usepackage[a4paper,margin=2.5cm]{geometry}
\usepackage[utf8]{inputenc}
\usepackage{amsmath,amssymb,amsthm,mathtools,bm}
\usepackage{microtype}
\usepackage{enumitem}
\usepackage{booktabs,array,multirow}
\usepackage{graphicx}
\usepackage{xcolor}
\usepackage{tikz}
\usetikzlibrary{arrows.meta,positioning,calc,decorations.pathmorphing,shapes.geometric,fit,backgrounds}
\usepackage{pgfplots}
\usepgfplotslibrary{groupplots}
\pgfplotsset{compat=1.17}
\usepackage{algorithm}
\usepackage{algpseudocode}
\usepackage{listings}
\usepackage[most]{tcolorbox}
\usepackage{fancyhdr}
\usepackage{float}
\usepackage[plainpages=false,pdfpagelabels,colorlinks=true,linkcolor=blue!50!black,citecolor=green!40!black,urlcolor=blue!50!black]{hyperref}
\usepackage[nameinlink,noabbrev]{cleveref}
% --- silence harmless warnings / make anchors unique ---
\makeatletter
\gdef\theHALG@line{\thealgorithm.\arabic{ALG@line}}
\makeatother
\hbadness=10000
\hfuzz=2pt
\vbadness=10000

\setlength{\parskip}{4pt}\emergencystretch=3em
\setlength{\parindent}{0pt}
\linespread{1.07}
\setcounter{tocdepth}{2}

% ---------- header / footer ----------
\pagestyle{fancy}
\fancyhf{}
\renewcommand{\chaptermark}[1]{\markboth{#1}{}}
\fancyhead[L]{\small\itshape\leftmark}
\fancyhead[R]{\small\itshape Elliptic Curve Cryptography}
\fancyfoot[C]{\thepage}
\renewcommand{\headrulewidth}{0.3pt}
\fancypagestyle{plain}{\fancyhf{}\fancyfoot[C]{\thepage}\renewcommand{\headrulewidth}{0pt}}

% ---------- theorem environments ----------
\theoremstyle{plain}
\newtheorem{theorem}{Theorem}[chapter]
\newtheorem{lemma}[theorem]{Lemma}
\newtheorem{proposition}[theorem]{Proposition}
\newtheorem{corollary}[theorem]{Corollary}
\newtheorem{conjecture}[theorem]{Conjecture}
\theoremstyle{definition}
\newtheorem{definition}[theorem]{Definition}
\newtheorem{example}[theorem]{Example}
\newtheorem{problem}[theorem]{Problem}
\theoremstyle{remark}
\newtheorem{remark}[theorem]{Remark}
\newtheorem*{notation}{Notation}
\newtheorem{exercise}{Exercise}[chapter]

\crefname{theorem}{Theorem}{Theorems}
\crefname{lemma}{Lemma}{Lemmas}
\crefname{proposition}{Proposition}{Propositions}
\crefname{corollary}{Corollary}{Corollaries}
\crefname{definition}{Definition}{Definitions}
\crefname{example}{Example}{Examples}
\crefname{remark}{Remark}{Remarks}
\crefname{exercise}{Exercise}{Exercises}
\crefname{figure}{Figure}{Figures}
\crefname{table}{Table}{Tables}
\crefname{algorithm}{Algorithm}{Algorithms}
\crefname{chapter}{Lecture}{Lectures}
\crefname{section}{Section}{Sections}

% ---------- coloured boxes ----------
\newtcolorbox{keyidea}[1][Key idea]{colback=blue!4!white,colframe=blue!55!black,
  fonttitle=\bfseries\small,title={#1},breakable,left=4pt,right=4pt,top=3pt,bottom=3pt}
\newtcolorbox{warningbox}[1][Warning]{colback=red!4!white,colframe=red!60!black,
  fonttitle=\bfseries\small,title={#1},breakable,left=4pt,right=4pt,top=3pt,bottom=3pt}
\newtcolorbox{goalsbox}[1][Learning goals]{colback=green!4!white,colframe=green!40!black,
  fonttitle=\bfseries\small,title={#1},breakable,left=4pt,right=4pt,top=3pt,bottom=3pt}

% ---------- macros ----------
\newcommand{\Z}{\mathbb{Z}}
\newcommand{\N}{\mathbb{N}}
\newcommand{\Q}{\mathbb{Q}}
\newcommand{\R}{\mathbb{R}}
\newcommand{\C}{\mathbb{C}}
\newcommand{\F}{\mathbb{F}}
\newcommand{\Fp}{\mathbb{F}_p}
\newcommand{\Fq}{\mathbb{F}_q}
\newcommand{\Fqbar}{\overline{\mathbb{F}}_q}
\renewcommand{\P}{\mathbb{P}}
\newcommand{\A}{\mathbb{A}}
\newcommand{\E}{E}
\newcommand{\Oinf}{\mathcal{O}}
\newcommand{\Kbar}{\overline{K}}
\newcommand{\ord}{\operatorname{ord}}
\newcommand{\End}{\operatorname{End}}
\newcommand{\Gal}{\operatorname{Gal}}
\newcommand{\Div}{\operatorname{Div}}
\newcommand{\Pic}{\operatorname{Pic}}
\newcommand{\divv}{\operatorname{div}}
\newcommand{\Tr}{\operatorname{Tr}}
\newcommand{\GL}{\operatorname{GL}}
\newcommand{\lcm}{\operatorname{lcm}}
\newcommand{\Frob}{\varphi}
\newcommand{\Hom}{\operatorname{Hom}}
\newcommand{\legendre}[2]{\left(\frac{#1}{#2}\right)}
\newcommand{\bigO}{\mathcal{O}}
\newcommand{\softO}{\widetilde{\mathcal{O}}}
\newcommand{\Mul}{\mathbf{M}}
\newcommand{\Sqr}{\mathbf{S}}
\newcommand{\Inv}{\mathbf{I}}
\newcommand{\ECDLP}{\mathsf{ECDLP}}
\newcommand{\CDH}{\mathsf{CDH}}
\newcommand{\DDH}{\mathsf{DDH}}
\newcommand{\ELsection}{\chapter}
\newcommand{\sample}{\xleftarrow{\$}}

% ---------- listings ----------
\lstdefinestyle{py}{language=Python,basicstyle=\ttfamily\footnotesize,keywordstyle=\color{blue!60!black}\bfseries,
  commentstyle=\color{green!40!black}\itshape,stringstyle=\color{red!50!black},showstringspaces=false,
  frame=single,framerule=0.3pt,rulecolor=\color{black!40},numbers=left,numberstyle=\tiny\color{black!50},
  breaklines=true,columns=fullflexible,xleftmargin=1.5em,framexleftmargin=1.5em,tabsize=4}
\lstdefinestyle{plain}{basicstyle=\ttfamily\footnotesize,frame=single,framerule=0.3pt,breaklines=true,columns=fullflexible}

\algrenewcommand\algorithmicrequire{\textbf{Input:}}
\algrenewcommand\algorithmicensure{\textbf{Output:}}

\title{\Huge\bfseries EC-ElGamal Encryption\\[0.3cm]\Large A self-contained introduction with worked examples}
\author{\large Companion handout to the lecture notes on Elliptic Curve Cryptography}
\date{\today}

\begin{document}
\maketitle
\mainmatter
\ELsection{EC-ElGamal encryption, step by step}\label{sec:ecelgamal}

\begin{keyidea}[What you need to know beforehand \,--\, only four facts]
\begin{enumerate}[nosep]
\item The points of an elliptic curve $E$ over $\F_p$ (plus a point at infinity $\Oinf$) form a group written \emph{additively}: we can add points $P+Q$, negate $-(x,y)=(x,-y)$, and subtract $P-Q=P+(-Q)$. The neutral element is $\Oinf$.
\item $G$ is a public base point whose multiples $G,2G,3G,\dots,nG=\Oinf$ form a cyclic group of (prime) order $n$. Computing $kG$ is fast (double-and-add). Finding $k$ from $kG$ is the hard problem (ECDLP).
\item Scalar multiples commute: $k(dG)=d(kG)=(kd)G$.
\item ECDH: Alice sends $A=aG$, Bob sends $B=bG$, and both compute the same point $aB=bA=abG$.
\end{enumerate}
That is all. If you remember ECDH, ElGamal will take ten minutes.
\end{keyidea}

\paragraph{The big idea: ElGamal = ECDH + a one-time pad.}
In ECDH the two parties end up with a shared secret point $abG$ and then turn it into a key. ElGamal does the same thing but with the roles arranged for \emph{encryption}:
\begin{itemize}[nosep]
\item Bob publishes a \emph{long-term} public key $Q=dG$ (it plays the role of ``$B=bG$'' in ECDH; $d$ is his private key).
\item Alice, who wants to send a message $M$, picks a fresh \emph{ephemeral} secret $k$ and sends $C_1=kG$ (the role of ``$A=aG$'').
\item Now Alice and Bob \emph{both} know the shared point $S=kQ=dC_1=kdG$, exactly as in ECDH.
\item Alice uses $S$ as a one-time mask: she \emph{adds} it to the message, $C_2=M+S$. Only someone who knows $S$ can undo this.
\end{itemize}
The ciphertext is the pair $(C_1,C_2)$. Because Alice's secret $k$ is fresh every time, the mask $S$ is different every time, like a fresh one-time pad.

\begin{figure}[ht]
\centering
\begin{tikzpicture}[>=Stealth]
\node[draw,rounded corners,fill=blue!8,minimum width=3.2cm] (A) at (0,0) {Alice (sender)};
\node[draw,rounded corners,fill=blue!8,minimum width=3.2cm] (B) at (9,0) {Bob (receiver)};
\draw[dashed,gray] (0,-0.4)--(0,-6.6); \draw[dashed,gray] (9,-0.4)--(9,-6.6);
\node[font=\small,anchor=east,align=right] at (8.7,-1.1) {private key $d$\\ public key $Q=dG$};
\draw[<-,thick] (0.1,-2.3)--(8.9,-2.3) node[midway,above,font=\small]{Bob publishes $Q$ (once)};
\node[font=\small,anchor=west,align=left] at (0.3,-3.3) {pick fresh random $k$\\ $C_1=kG$\\ mask $S=kQ$\\ $C_2=M+S$};
\draw[->,very thick,green!50!black] (0.1,-4.7)--(8.9,-4.7) node[midway,above,font=\small]{send ciphertext $(C_1,C_2)$};
\node[font=\small,anchor=east,align=right] at (8.7,-5.7) {mask $S=dC_1$ (equals $kQ$!)\\ $M=C_2-S$};
\node[draw,dotted,fill=red!6,font=\small,align=center,anchor=west] at (0.3,-5.9) {Eve sees $G,\,Q,\,C_1,\,C_2$\\ but not $S$};
\end{tikzpicture}
\caption{EC-ElGamal as ECDH plus masking. The point $C_1=kG$ is Alice's ECDH public value; the shared point $S=kQ=dC_1$ is the ECDH secret; it is used to hide $M$ by adding it.}
\label{fig:EGoverview}
\end{figure}

\paragraph{Dictionary: from classical ElGamal to the elliptic-curve version.}
If you know ElGamal in $\Z_p^\times$ (as in Paar's book), the elliptic-curve version is the same scheme with the group written additively: \emph{multiplication becomes addition} and \emph{exponentiation becomes scalar multiplication}.
\begin{center}
\small
\renewcommand{\arraystretch}{1.25}
\begin{tabular}{@{}lll@{}}
\toprule
 & Classical ElGamal in $\Z_p^\times$ & EC-ElGamal on $E(\F_p)$\\
\midrule
public base & $g$ & $G$\\
private key & $d$ & $d$\\
public key & $\beta=g^d$ & $Q=dG$\\
message & integer $m\in\Z_p^\times$ & a \emph{point} $M\in\langle G\rangle$\\
ephemeral secret & $k$ & $k$\\
first part & $c_1=g^k$ & $C_1=kG$\\
mask & $\beta^k$ & $kQ$\\
second part & $c_2=m\cdot\beta^k$ & $C_2=M+kQ$\\
decryption & $m=c_2\cdot(c_1^{\,d})^{-1}$ & $M=C_2-dC_1$\\
\bottomrule
\end{tabular}
\end{center}

\subsubsection*{The scheme}

\begin{definition}[EC-ElGamal]\label{def:ECEG}
Public parameters: curve $E$ over $\F_p$, point $G$ of prime order $n$.
\begin{enumerate}[nosep,label=\textbf{\arabic*.}]
\item \textbf{Key generation (Bob).} Choose $d\sample\{1,\dots,n-1\}$ uniformly. Public key $Q=dG$; private key $d$.
\item \textbf{Encryption (Alice)} of a message point $M\in\langle G\rangle$: choose a fresh random $k\sample\{1,\dots,n-1\}$ and output $\big(C_1,C_2\big)=\big(kG,\ M+kQ\big)$.
\item \textbf{Decryption (Bob)}: compute $M=C_2-dC_1$.
\end{enumerate}
\end{definition}

\begin{theorem}[Correctness]
Decryption returns $M$.
\end{theorem}
\begin{proof}
$dC_1=d(kG)=k(dG)=kQ$, hence $C_2-dC_1=(M+kQ)-kQ=M$.
\end{proof}

\paragraph{Cost.} Encryption needs two scalar multiplications ($kG$ and $kQ$) and one point addition; decryption needs one scalar multiplication and one subtraction. The ciphertext consists of \emph{two} points, i.e.\ it is about twice as long as the message point.

\subsubsection*{A problem you must solve first: messages are not points}

The scheme encrypts a \emph{point} $M$, but we want to encrypt numbers, text, keys. The classical (Koblitz) fix is to embed a small integer $m$ into a point:

\begin{keyidea}[Koblitz embedding with parameter $\kappa$]
To encode $m\in\{0,1,\dots,\lfloor p/\kappa\rfloor-1\}$: try $x=\kappa m,\ \kappa m+1,\ \dots,\ \kappa m+\kappa-1$ and take the first $x$ for which $x^3+ax+b$ is a square modulo $p$ (test with Euler's criterion); let $y$ be a square root (choose the smaller one). Then $M_m=(x,y)$. To decode, output $m=\lfloor x/\kappa\rfloor$ (only the $x$-coordinate is needed). Each trial succeeds with probability about $\frac12$, so encoding fails with probability about $2^{-\kappa}$; real systems take $\kappa\approx30$--$60$.
\end{keyidea}

We will work on a curve small enough to do everything by hand.

\paragraph{The toy curve.} $E:\ y^2=x^3+2x+2$ over $\F_{17}$, with base point $G=(5,1)$ of prime order $n=19$. All multiples of $G$ (you will use this table all the time):
\begin{center}
\scriptsize\setlength{\tabcolsep}{2.6pt}
\begin{tabular}{@{}c|*{9}{c}@{}}
$k$ & 1&2&3&4&5&6&7&8&9\\\hline
$kG$ & $(5,1)$&$(6,3)$&$(10,6)$&$(3,1)$&$(9,16)$&$(16,13)$&$(0,6)$&$(13,7)$&$(7,6)$\\[4pt]
$k$ & 10&11&12&13&14&15&16&17&18\\\hline
$kG$ & $(7,11)$&$(13,10)$&$(0,11)$&$(16,4)$&$(9,1)$&$(3,16)$&$(10,11)$&$(6,14)$&$(5,16)$
\end{tabular}
\end{center}
Remember $-kG=(19-k)G$, i.e.\ negation flips $y\mapsto-y$; for instance $-(6,3)=(6,14)$. And $19G=\Oinf$.

\paragraph{Embedding table} for $\kappa=3$, so $m\in\{0,\dots,4\}$ (the squares modulo $17$ are $1,2,4,8,9,13,15,16$; one checks that $f(x)=x^3+2x+2\bmod17$ is a square exactly for $x=0,3,5,6,7,9,10,13,16$):
\begin{center}
\small
\begin{tabular}{@{}c|ccccc@{}}
$m$ & 0 & 1 & 2 & 3 & 4\\\hline
$x$ tried & $0$ \checkmark & $3$ \checkmark & $6$ \checkmark & $9$ \checkmark & $12$ (no), $13$ \checkmark\\
$M_m$ & $(0,6)$ & $(3,1)$ & $(6,3)$ & $(9,1)$ & $(13,7)$\\
$=$ & $7G$ & $4G$ & $2G$ & $14G$ & $8G$
\end{tabular}
\end{center}

\subsubsection*{Worked example 1: encrypt and decrypt}

\begin{example}[Bob's key pair]\label{ex:EGkeys}
Bob chooses the private key $d=7$ and computes $Q=7G=(0,6)$ (table). He publishes $Q$.
\end{example}

\begin{example}[Alice encrypts the message $m=4$]\label{ex:EGenc}
Alice wants to send $m=4$. She embeds it: $M=(13,7)=8G$. She chooses the fresh random nonce $k=3$ and computes
\begin{align*}
C_1&=kG=3G=(10,6),\\
S&=kQ=3\cdot(0,6)=3\cdot7G=21G=2G=(6,3)\qquad\text{(the mask)},\\
C_2&=M+S=8G+2G=10G=(7,11).
\end{align*}
She sends $(C_1,C_2)=\big((10,6),(7,11)\big)$.
\end{example}

\begin{example}[Bob decrypts]\label{ex:EGdec}
Bob computes the mask from $C_1$ with his private key:
$S'=dC_1=7\cdot(10,6)=7\cdot3G=21G=2G=(6,3)$. This is the \emph{same} point $S$ that Alice used. He subtracts it:
\[
M=C_2-S'=(7,11)+\big(-(6,3)\big)=(7,11)+(6,14).
\]
The slope is $\lambda=\frac{14-11}{6-7}=\frac{3}{-1}\equiv14$, so $x_3=14^2-7-6=183\equiv13$ and $y_3=14(7-13)-11=-95\equiv7$, giving $M=(13,7)$. Decoding: $m=\lfloor13/3\rfloor=4$. \checkmark
\end{example}

\paragraph{What does Eve see?} Only $G=(5,1)$, $Q=(0,6)$, $C_1=(10,6)$, $C_2=(7,11)$. To recover $M$ she needs the mask $S=kQ=(6,3)$. Computing it needs $k$ (from $C_1=kG$, an ECDLP) or $d$ (from $Q=dG$, an ECDLP). On this toy curve she simply reads $k=3$ off the table; on a $256$-bit curve that is infeasible. Computing $S$ from $(G,Q,C_1)$ without either secret is exactly the Diffie--Hellman problem $\CDH$.

\begin{example}[Same message, new ciphertext]
Alice encrypts $m=4$ again with a new nonce $k=11$: $C_1=11G=(13,10)$, $S=11Q=77G=G=(5,1)$ (as $77=4\cdot19+1$), $C_2=8G+G=9G=(7,6)$. The ciphertext $\big((13,10),(7,6)\big)$ looks completely different, yet Bob decrypts it: $dC_1=7\cdot11G=77G=G$ and $9G-G=8G=(13,7)$. \checkmark
This randomisation is essential: if encryption were deterministic, anyone could encrypt guesses and compare.
\end{example}

\subsubsection*{Why is it secure?}

\begin{theorem}[Semantic security]
If the Decisional Diffie--Hellman problem is hard in $\langle G\rangle$, EC-ElGamal is IND-CPA secure: an adversary who sees the public key and a ciphertext learns nothing about which of two chosen messages was encrypted.
\end{theorem}
\begin{proof}[Idea]
The adversary sees $Q=dG$, $C_1=kG$ and $C_2=M+S$. If $S=kdG$ were replaced by a uniformly random point $R$, then $C_2=M+R$ would be uniformly distributed whatever $M$ is, exactly as with a one-time pad, and would carry no information about $M$. DDH hardness says that the genuine $S=kdG$ cannot be distinguished from such a random $R$ given $(G,dG,kG)$. So the real ciphertext is indistinguishable from one that hides $M$ perfectly.
\end{proof}

\paragraph{Intuition.} The point $S$ is a ``fresh one-time pad'' which only Alice (who knows $k$) and Bob (who knows $d$) can compute. This is the same fact that makes ECDH secure.

\subsubsection*{Three pitfalls (and what they teach)}

\paragraph{Pitfall 1: never reuse the nonce $k$.}
If two messages are encrypted with the same $k$, the same mask $S$ is added twice, like reusing a one-time pad.
\begin{example}[Nonce reuse]
Alice encrypts $M_1=(13,7)=8G$ and $M_2=(6,3)=2G$, both with $k=3$. The mask is $S=2G$ in both cases, so $C_1$ is the same and
$C_2=8G+2G=10G=(7,11)$, $C_2'=2G+2G=4G=(3,1)$.
Eve notices the repeated $C_1=(10,6)$ and computes $C_2-C_2'=(7,11)-(3,1)=6G=(16,13)=M_1-M_2$. If she also knows (or guesses) one plaintext, say $M_1$, she gets the mask $S=C_2-M_1=(6,3)$ and then $M_2=C_2'-S=(3,1)-(6,3)=2G=(6,3)$. All without the private key.
\end{example}
The nonce must be uniformly random and fresh for \emph{every} encryption (a bad random-number generator is as bad as reuse).

\paragraph{Pitfall 2: ciphertexts can be modified (malleability).}
\begin{example}[An attacker edits the plaintext]
Take Alice's ciphertext $\big((10,6),(7,11)\big)$ of $m=4$ from \cref{ex:EGenc}. Mallory replaces $C_2$ by $C_2+G=(7,11)+(5,1)=11G=(13,10)$ and forwards $\big((10,6),(13,10)\big)$. Bob decrypts $11G-2G=9G=(7,6)$ and decodes $m=\lfloor7/3\rfloor=2$. The message changed from $4$ to $2$, and Bob notices nothing, although Mallory never learned the key.
\end{example}
In general adding a point $T$ to $C_2$ changes the plaintext to $M+T$. This is why EC-ElGamal alone is not secure against \emph{active} attackers (it is not IND-CCA2). Practical systems use ECIES (Lecture~4), which authenticates the ciphertext.

\paragraph{Pitfall 3 (a feature in disguise): the scheme is additively homomorphic.}
\begin{example}[Adding encrypted messages]
Encrypt $M_1=(13,7)=8G$ with $k=3$: $(C_1,C_2)=\big((10,6),(7,11)\big)$. Encrypt $M_2=(3,1)=4G$ (this is $m=1$) with $k=11$: $(C_1',C_2')=\big((13,10),(9,16)\big)$. Add the ciphertexts component-wise:
\[
\big(C_1+C_1',\ C_2+C_2'\big)=\big((10,6)+(13,10),\ (7,11)+(9,16)\big)=\big((9,1),\ (3,16)\big).
\]
Bob decrypts: $d(9,1)=7\cdot14G=98G=3G=(10,6)$, and $(3,16)-(10,6)=15G-3G=12G=(0,11)$. And indeed $M_1+M_2=8G+4G=12G=(0,11)$. \checkmark\\
So one can add plaintexts \emph{without decrypting}. This is used in electronic voting: encode a vote $v\in\{0,1\}$ as the point $vG$ (``exponential ElGamal''), add all encrypted ballots, decrypt once to get $(\sum v)G$ and recover the tally $\sum v$ by trying $0G,1G,2G,\dots$ (cheap for small tallies).
\end{example}

\paragraph{Practical remarks.}
\begin{itemize}[nosep]
\item Before using a received ciphertext, check that $C_1,C_2$ are valid points of the curve in the correct subgroup (otherwise invalid-curve attacks, Lecture 4); reject $C_1=\Oinf$.
\item The message $M$ must lie in the subgroup $\langle G\rangle$; on prime-order curves that is every point.
\item The embedded message is short (about $250$ bits on a $256$-bit curve). Real systems encrypt long data with a symmetric cipher and use ElGamal only to transport the key: this is \textbf{ECIES}: $C_1=kG$, key $=\mathrm{KDF}(x(kQ))$, then authenticated encryption of the data. ECIES avoids both the embedding problem and malleability.
\end{itemize}

\subsubsection*{Python implementation (toy curve)}

The following code reproduces every number in this section (it uses the \texttt{Curve} class of the appendix with $E(\F_{17})$). It is for learning only.
\begin{lstlisting}[style=py]
import secrets
from ecc_toy import Curve                 # class Curve from the appendix

E = Curve(17, 2, 2, G=(5, 1), n=19)       # y^2 = x^3 + 2x + 2 over F_17, order of G = 19
KAPPA = 3

def keygen():
    d = secrets.randbelow(E.n - 1) + 1    # private key
    return d, E.mul(d, E.G)               # (d, Q = dG)

def embed(m, kappa=KAPPA):                # integer -> curve point (Koblitz)
    for j in range(kappa):
        x = kappa * m + j
        rhs = (x**3 + E.a * x + E.b) % E.p
        ys = [y for y in range(E.p) if y * y % E.p == rhs]
        if ys: return (x, min(ys))
    raise ValueError("embedding failed")

def unembed(P, kappa=KAPPA):              # curve point -> integer
    return P[0] // kappa

def encrypt(Q, M, k=None):
    k = k or secrets.randbelow(E.n - 1) + 1          # FRESH random nonce
    return E.mul(k, E.G), E.add(M, E.mul(k, Q))      # (C1, C2) = (kG, M + kQ)

def decrypt(d, C):
    C1, C2 = C
    return E.add(C2, E.neg(E.mul(d, C1)))            # M = C2 - d*C1

d, Q = 7, E.mul(7, E.G)
M = embed(4)                                         # (13, 7)
C = encrypt(Q, M, k=3)                               # ((10, 6), (7, 11))
print(C, decrypt(d, C), unembed(decrypt(d, C)))      # ... (13, 7) 4
\end{lstlisting}

\begin{keyidea}[One-page summary]
\begin{itemize}[nosep]
\item \textbf{Keys:} private $d$; public $Q=dG$.
\item \textbf{Encrypt $M$:} pick random $k$; send $\big(kG,\ M+kQ\big)$.
\item \textbf{Decrypt:} $M=C_2-dC_1$, because $dC_1=kQ$ (the ECDH secret).
\item \textbf{Security:} $kQ$ is a fresh one-time pad; secure if DDH is hard.
\item \textbf{Must do:} fresh random $k$ every time; embed messages as points; validate points.
\item \textbf{Limits:} malleable (edit $C_2$ edits $M$), ciphertext twice the size, short messages $\Rightarrow$ use ECIES in practice.
\item \textbf{Bonus:} additively homomorphic $\Rightarrow$ e-voting, private sums.
\end{itemize}
\end{keyidea}

\subsubsection*{Exercises (with answers)}

\begin{exercise}
With Bob's key $d=7$, $Q=(0,6)$ encrypt $M=(9,1)$ (this is $m=3$) using $k=2$, then decrypt. \emph{Answer:} $C_1=2G=(6,3)$, $S=2Q=14G=(9,1)$, $C_2=14G+14G=28G=9G=(7,6)$; ciphertext $\big((6,3),(7,6)\big)$. Decryption: $dC_1=14G$ and $9G-14G=-5G=14G=(9,1)$. 
\end{exercise}

\begin{exercise}
Decrypt $\big((13,10),(0,11)\big)$ with $d=7$ and decode the message. \emph{Answer:} $C_1=11G$, $dC_1=77G=G$, $M=12G-G=11G=(13,10)$, so $m=\lfloor13/3\rfloor=4$ (decoding only uses the $x$-coordinate).
\end{exercise}

\begin{exercise}
What happens if Alice chooses $k=19$ (or $k=0$)? \emph{Answer:} $C_1=\Oinf$ and $S=\Oinf$, so $C_2=M$: the plaintext is sent in the clear. This is why $k$ is drawn from $\{1,\dots,n-1\}$.
\end{exercise}

\begin{exercise}
Two ciphertexts $\big((10,6),(7,11)\big)$ and $\big((10,6),(3,1)\big)$ are intercepted, and Eve knows the first plaintext is $(13,7)$. Recover the second. \emph{Answer:} the same $C_1$ means the same $k$; $S=(7,11)-(13,7)=(6,3)$; $M_2=(3,1)-(6,3)=(6,3)$, i.e.\ $m=2$.
\end{exercise}

\begin{exercise}
In an election with encrypted votes $v_i\in\{0,1\}$ encoded as $v_iG$, explain how to obtain the tally without decrypting individual votes, and why it only works for small tallies. \emph{Answer:} add all ciphertexts (component-wise), decrypt the sum once to get $TG$, $T=\sum v_i$; find $T$ by testing $0G,1G,\dots$ (a discrete logarithm of a small number: cheap only if $T$ is at most about $2^{40}$).
\end{exercise}

\begin{exercise}
Show that $(C_1,\,2C_2)$ does \emph{not} decrypt to $2M$, but $(2C_1,\,2C_2)$ does. \emph{Answer:} decrypting $(C_1,2C_2)$ gives $2C_2-dC_1=2M+kQ\ne2M$, whereas $(2C_1,2C_2)$ gives $2C_2-2dC_1=2M+2kQ-2kQ=2M$.
\end{exercise}

\end{document}
